Una porta s'obre i es tanca mitjançant un dispositiu fotoelèctric. La longitud d'ona de la radiació electromagnètica utilitzada és de 850~nm i l'energia mínima d'extracció del material fotodetector és d'1,20~eV. Calculeu:

\begin{apartat}{1}
L'energia cinètica dels fotoelectrons emesos i la longitud d'ona de De Broglie associada a aquests electrons.
\end{apartat}

\begin{apartat}{1}
La longitud d'ona que hauria de tenir una radiació electromagnètica incident per a duplicar l'energia cinètica dels fotoelectrons emesos de l'apartat \textit{a}.
\end{apartat}

\dades{%
  $m_\text{electró} = 9,11\times 10^{-31}\ \mathrm{kg}$\\
  $Q_\text{electró} = -1,60\times 10^{-19}\ \mathrm{C}$\\
  $h = 6,63\times 10^{-34}\ \mathrm{J\,s}$\\
  $c = 3,00\times 10^{8}\ \mathrm{m\,s^{-1}}$}
